\section{Introducción}

El presente informe desarrolla la documentación técnica de una base de datos para la gestión de procesos odontológicos. El dominio considerado incluye registro de pacientes, asignación de carpetas clínicas, programación de citas, atención mediante odontogramas, tratamientos, planes de pago, registro de pagos, compras de suministros e inventario.

El trabajo toma como referencia documental una tesis previa sobre el desarrollo de un sistema de información web para el consultorio odontológico LALYSDENT. En dicha tesis se documenta un sistema web implementado con Laravel, MySQL y arquitectura MVC. A diferencia de esa investigación, el presente informe no se enfoca en el frontend ni en la aplicación web, sino en el diseño, transformación, implementación y validación de la base de datos.

La propuesta fue desarrollada en etapas. Primero, se definieron los requisitos y supuestos del dominio. Luego, se elaboró un modelo conceptual en notación Chen usando diagrams.net. Posteriormente, se realizó la transformación al modelo relacional y se construyó un modelo lógico Crow's Foot en ER/Studio 2003. Finalmente, se implementó el diseño físico en Oracle Database 19c, ejecutado en Docker y administrado mediante DBeaver.

El documento incluye también un diccionario de datos completo, validaciones ejecutadas sobre Oracle 19c y una comparación técnica entre el enfoque MySQL/Laravel documentado en la tesis base y la transformación realizada hacia Oracle 19c.
